\section{Introduction and statement of the theorem}\label{sec:intro}
Let $m\in\Z_{\ge3}$. Let $K(m)$ be the quiver with two vertices and $m$ arrows from the first
vertex to the second. All quiver representations are over $\C$.
Let $\chi$ denote the ordinary topological Euler characteristic.
For positive coprime integers $a$ and $b$, let $\cM^s_{(a,b)}(K(m))$ denote the moduli space of stable representations of $K(m)$ of dimension $(a,b)$, where $a$ is the dimension at the source and $b$ at the target.
Let
\[
 E_m(a,b)=\chi\bigl(\cM^s_{(a,b)}(K(m))\bigr).
\]
Our stability requires
\begin{equation}\label{eq:stability}
 av-bu>0
\end{equation}
for every nonzero proper subrepresentation of dimension $(u,v)$.
For
\[
 \lambda_\pm=\frac{m\pm\sqrt{m^2-4}}2,
 \qquad I_m=[\lambda_-,\lambda_+],
\]
let
\[
 \cA_m=\{(a,b)\in\Z_{>0}^2:\gcd(a,b)=1,\
       \cM^s_{(a,b)}(K(m))\ne\emptyset\}.
\]

We prove Douglas' conjecture in the following
\cite[Conjecture 6.1]{Weist}
(see also~\cite{Okada}).
The study of black hole entropies motivates Douglas' conjecture~\cite{CS,Denef,DM,DFR,MPS,Weist}.

\begin{theorem}\label{thm:main}
For each integer $m\geq3$, there is a continuous function
$F_m:I_m\to\R_{\ge0}$ such that
\[
 F_m(\lambda_-)=F_m(\lambda_+)=0,\qquad
 0<F_m(r)<\infty\quad(\lambda_-<r<\lambda_+)
\]
and
\begin{equation}\label{eq:main-limit}
 \lim_{\substack{a\to\infty,\ (a,b)\in\cA_m\\\frac{b}{a}\to r\in I_m}}
 \frac{\log E_m(a,b)}a=F_m(r).
\end{equation}
\end{theorem}

By~\cite[Theorem 4.1]{Sch01} and~\cite[Theorem 6.1]{Sch92}, we have
\begin{equation}\label{eq:nonempty-domain}
 \cA_m=\{(a,b)\in\Z_{>0}^2:
       \gcd(a,b)=1,\ a^2+b^2-mab\le1\}.
\end{equation}
For $a^2+b^2-mab\le0$, see also~\cite[Theorems 2.7(1) and 4.1(1)]{Rei}.
Furthermore,~\eqref{eq:main-limit} is equivalent to the following: for every $r\in I_m$ and $\varepsilon>0$, there are $\delta>0$ and $A>0$
such that
\[
 \left|\frac{\log E_m(a,b)}a-F_m(r)\right|<\varepsilon
 \quad\text{if }(a,b)\in\cA_m,\ a>A,
       \left|\frac ba-r\right|<\delta.
\]

\subsection{The counting model and the comparison argument}
We use a counting formula for $mE_m(a,b)$~\cite{B} with quiver representations.
The finite model counts colored gradings on paths.
We construct maps between these sets of gradings and obtain a product estimate for their cardinalities.
The product estimate gives an exponential growth rate under repeated concatenation.

We first prove the limit in~\eqref{eq:main-limit} for rational $r$.
We then obtain the limit for irrational $r$ by comparison with rational values.
Finally, we prove continuity of $F_m$ and the endpoint values
$F_m(\lambda_-)=F_m(\lambda_+)=0$.

